\documentclass[../../../include/open-logic-section]{subfiles}

\begin{document}

\olfileid[id]{his}{bio}{noe}

\olsection{Emmy Noether}

\olphoto{noether-emmy}{Emmy Noether}

Emmy Noether (\textsc{ner}-ter) lahir di Erlangen, Jerman, pada 23 Maret
1882 dalam keluarga akademis kelas menengah atas. Dipuji sebagai ``ibu
aljabar modern,'' Noether memberikan sumbangan terobosan bagi matematika
dan fisika, meskipun menghadapi hambatan besar terhadap pendidikan
perempuan. Di Jerman pada masa itu, anak perempuan diharapkan memperoleh
pendidikan dalam bidang seni dan tidak diizinkan bersekolah di sekolah
persiapan perguruan tinggi. Namun, setelah mengikuti kuliah sebagai
pendengar di Universitas G\"{o}ttingen dan Universitas Erlangen (tempat
ayahnya menjadi profesor matematika), Noether akhirnya dapat mendaftar
sebagai mahasiswa di Erlangen pada tahun 1904, ketika kebijakan universitas
diubah untuk menerima mahasiswi. Ia meraih gelar doktor dalam matematika
pada tahun 1907.

Meskipun memenuhi syarat, Noether menghadapi banyak penolakan sepanjang
kariernya. Dari tahun 1908--1915, ia mengajar di Erlangen tanpa dibayar.
Pada masa ini, ia menarik perhatian David Hilbert, salah seorang
matematikawan terkemuka dunia ketika itu, yang mengundangnya ke
G\"{o}ttingen. Namun, perempuan dilarang memperoleh jabatan profesor, dan ia
hanya dapat memberikan kuliah atas nama Hilbert, sekali lagi tanpa dibayar.
Pada masa inilah ia membuktikan hasil yang sekarang dikenal sebagai teorema
Noether, yang hingga kini masih digunakan dalam fisika teoretis. Noether
akhirnya memperoleh hak mengajar pada tahun 1919. Tanggapan Hilbert terhadap
penolakan yang terus berlanjut dari rekan-rekan universitasnya konon adalah:
``Tuan-tuan, senat fakultas bukanlah tempat pemandian.''

Pada akhir dasawarsa 1920-an, ia memusatkan perhatian pada aljabar abstrak,
dan sumbangannya merevolusi bidang tersebut. Dalam pembuktiannya ia sering
menggunakan apa yang disebut syarat rantai menaik, yang menyatakan bahwa
tidak ada rantai tak hingga yang menaik ketat dari himpunan-himpunan
tertentu. Sebagai contoh, struktur aljabar tertentu yang sekarang dikenal
sebagai gelanggang Noetherian mempunyai sifat bahwa tidak ada barisan ideal
tak hingga $I_1 \subsetneq I_2 \subsetneq \dots$. Syarat itu dapat
digeneralisasi ke sebarang urutan parsial (dalam aljabar, syarat ini
menyangkut kasus khusus ideal yang diurutkan oleh relasi himpunan bagian),
dan kita juga dapat mempertimbangkan syarat rantai menurun yang dual, yakni
setiap barisan yang \emph{menurun} ketat dalam suatu urutan parsial pada
akhirnya berakhir. Jika suatu urutan parsial memenuhi syarat rantai menurun,
kita dapat menggunakan induksi sepanjang urutan itu dengan cara yang serupa
dengan penggunaan induksi sepanjang urutan $<$ pada~$\Nat$. Urutan demikian
disebut \emph{berlandasan baik} atau \emph{Noetherian}, dan prinsip bukti
yang bersesuaian disebut \emph{induksi Noetherian}.

Noether adalah seorang Yahudi, dan ketika Nazi berkuasa pada tahun 1933, ia
diberhentikan dari jabatannya. Beruntung, Noether dapat beremigrasi ke
Amerika Serikat untuk menduduki jabatan sementara di Bryn Mawr,
Pennsylvania. Selama berada di sana ia juga memberi kuliah di Princeton,
meskipun ia merasa universitas tersebut tidak ramah terhadap perempuan
\citep[81]{Dick1981}. Pada tahun 1935, Noether menjalani operasi untuk
mengangkat tumor rahim. Ia meninggal akibat infeksi yang timbul dari operasi
tersebut dan dimakamkan di Bryn Mawr.

\begin{reading}
Untuk biografi Noether, lihat \citet{Dick1981}. Perimeter Institute for
Theoretical Physics menyediakan secara daring kuliah-kuliah mereka mengenai
kehidupan dan pengaruh Noether \citep{Perimeter2015}. Jika Anda lelah
membaca, \emph{Stuff You Missed in History Class} mempunyai siniar mengenai
kehidupan dan pengaruh Noether \citep{Frey2015}. Kumpulan karya Noether
tersedia dalam bahasa Jerman aslinya \citep{Noether1983}.
\end{reading}

\end{document}
